% ============================================================================
% BOLUM 5: GELISTIRILEN SISTEM
% ============================================================================
% Kurumsal Temizlik Takip ve Yonetim Sistemi (KTTYS)
% Muhammet Anil YAGIZ - 213255046
% Kirikkale Universitesi - Bilgisayar Muhendisligi Bolumu
% ============================================================================

\chapter{GELISTIRILEN SISTEM}
\label{chap:sistem-gelistirme}

% ============================================================================
% 5.1 GIRIS
% ============================================================================
\section{Giris}
\label{sec:sistem-giris}

Bu bolumde, Kurumsal Temizlik Takip ve Yonetim Sisteminin (KTTYS) 
teknik detaylari, veritabani tasarimi, API yapisi, kullanici arayuzu 
ve is akislari detayli olarak aciklanmaktadir.

% ============================================================================
% 5.2 SISTEM MIMARISI
% ============================================================================
\section{Sistem Mimarisi}
\label{sec:sistem-mimarisi}

KTTYS, uc katmanli bir mimari uzerine insa edilmistir: Sunum 
Katmani (Presentation Layer), Is Mantigi Katmani (Business Logic 
Layer) ve Veri Katmani (Data Layer).

\begin{figure}[H]
\centering
\begin{tikzpicture}[
    node distance=1.5cm,
    layer/.style={rectangle, draw, fill=blue!20, minimum width=11cm, 
                  minimum height=1.5cm, text centered, font=\bfseries},
    arrow/.style={->, thick, >=stealth}
]
    % Katmanlar
    \node[layer, fill=green!30] (presentation) {Sunum Katmani (HTML5, CSS3, JavaScript, PWA)};
    \node[layer, fill=orange!30, below=of presentation] (business) 
        {Is Mantigi Katmani (Node.js, Express.js, REST API)};
    \node[layer, fill=purple!30, below=of business] (data) 
        {Veri Katmani (SQLite, sql.js)};
    
    % Oklar
    \draw[arrow, <->] (presentation) -- (business) 
        node[midway, right, font=\small] {HTTP/REST};
    \draw[arrow, <->] (business) -- (data) 
        node[midway, right, font=\small] {SQL};
        
    % Kullanici
    \node[above=0.8cm of presentation] (user) {Kullanici};
    \draw[arrow] (user) -- (presentation);
    
\end{tikzpicture}
\caption{KTTYS Uc Katmanli Sistem Mimarisi}
\label{fig:sistem-mimarisi}
\end{figure}

\subsection{Sunum Katmani}
\label{subsec:sunum-katmani}

Sunum katmani, kullanicilarin sistemle etkilesime girdigi web 
arayuzunu icermektedir. Single Page Application (SPA) mimarisi 
benimsenmis olup, sayfa yenilemesi olmadan dinamik icerik 
guncellemesi saglanmaktadir.

Sunum katmaninin bilesenleri:

\begin{itemize}
    \item \textbf{index.html:} Ana HTML dosyasi, uygulama iskeleti
    \item \textbf{style.css:} Responsive CSS stilleri
    \item \textbf{app.js:} Uygulama mantigi ve UI guncellemeleri
    \item \textbf{api.js:} Backend API cagri fonksiyonlari
    \item \textbf{sw.js:} Service Worker, PWA destegi
    \item \textbf{manifest.json:} PWA manifest dosyasi
\end{itemize}

\subsection{Is Mantigi Katmani}
\label{subsec:is-mantigi-katmani}

Is mantigi katmani, Express.js framework'u uzerinde calisan 
RESTful API servislerini icermektedir. Bu katman, kullanici 
isteklerini isler, is kurallarini uygular ve veritabani 
islemlerini koordine eder.

\subsection{Veri Katmani}
\label{subsec:veri-katmani}

Veri katmani, SQLite veritabanini ve sql.js kutuphanesini 
icermektedir. Tum veriler tek bir .db dosyasinda 
saklanmaktadir.

% ============================================================================
% 5.3 VERITABANI TASARIMI
% ============================================================================
\section{Veritabani Tasarimi}
\label{sec:veritabani-tasarimi}

KTTYS veritabani, iliskisel veritabani modeli uzerine tasarlanmis 
olup, yedi adet tablo icermektedir.

\subsection{Varlik Iliski Diyagrami (ER Diagram)}
\label{subsec:er-diyagram}

Veritabani semasinin varlik-iliski diyagrami Sekil 
\ref{fig:er-diagram}'da gosterilmektedir.

\begin{figure}[H]
\centering
\begin{tikzpicture}[
    node distance=2cm,
    entity/.style={rectangle, draw=black, fill=blue!20, 
                   minimum width=2.5cm, minimum height=0.8cm,
                   text centered, font=\bfseries\small}
]
    % Varliklar
    \node[entity] (faculties) {FACULTIES};
    \node[entity, right=2.5cm of faculties] (users) {USERS};
    \node[entity, below=of faculties] (buildings) {BUILDINGS};
    \node[entity, below=of buildings] (departments) {DEPARTMENTS};
    \node[entity, right=2.5cm of departments] (locations) {LOCATIONS};
    \node[entity, right=2.5cm of users] (schedules) {SCHEDULES};
    \node[entity, below=of schedules] (tasks) {TASKS};
    \node[entity, below=of locations] (holidays) {HOLIDAYS};
    
    % Iliskiler
    \draw[->] (faculties) -- (users) node[midway, above] {1:N};
    \draw[->] (faculties) -- (buildings) node[midway, left] {1:N};
    \draw[->] (buildings) -- (departments) node[midway, left] {1:N};
    \draw[->] (departments) -- (locations) node[midway, above] {1:N};
    \draw[->] (locations) -- (schedules);
    \draw[->] (users) -- (schedules);
    \draw[->] (users) -- (tasks);
    \draw[->] (locations) -- (tasks);
    
\end{tikzpicture}
\caption{KTTYS Varlik Iliski Diyagrami}
\label{fig:er-diagram}
\end{figure}

Sekil \ref{fig:er-diagram}'daki ER diyagrami, sistemdeki temel 
varliklar ve aralarindaki iliskileri gostermektedir. Fakulteler 
birden cok kullanici ve binaya sahip olabilmektedir. Binalar 
departmanlari, departmanlar ise lokasyonlari icermektedir.

\subsection{Veritabani Tablolari}
\label{subsec:veritabani-tablolari}

\subsubsection{Faculties (Fakulteler) Tablosu}
\label{subsubsec:faculties}

Fakulteler tablosu, universitedeki fakultelerin bilgilerini 
saklamaktadir.

\begin{lstlisting}[language=SQL, caption={Faculties Tablosu Semasi}]
CREATE TABLE IF NOT EXISTS faculties (
    id INTEGER PRIMARY KEY AUTOINCREMENT,
    name TEXT NOT NULL,
    code TEXT UNIQUE NOT NULL,
    color TEXT DEFAULT '#1e40af',
    is_active INTEGER DEFAULT 1,
    created_at DATETIME DEFAULT CURRENT_TIMESTAMP
);
\end{lstlisting}

\subsubsection{Users (Kullanicilar) Tablosu}
\label{subsubsec:users}

Kullanicilar tablosu, sistemdeki tum kullanicilarin bilgilerini 
ve rollerini saklamaktadir.

\begin{lstlisting}[language=SQL, caption={Users Tablosu Semasi}]
CREATE TABLE IF NOT EXISTS users (
    id INTEGER PRIMARY KEY AUTOINCREMENT,
    faculty_id INTEGER REFERENCES faculties(id),
    username TEXT UNIQUE NOT NULL,
    password_hash TEXT NOT NULL,
    full_name TEXT NOT NULL,
    role TEXT CHECK(role IN ('super_admin', 'faculty_admin', 
                              'supervisor', 'staff')) NOT NULL,
    is_active INTEGER DEFAULT 1,
    created_at DATETIME DEFAULT CURRENT_TIMESTAMP
);
\end{lstlisting}

\subsubsection{Tasks (Gorevler) Tablosu}
\label{subsubsec:tasks}

\begin{lstlisting}[language=SQL, caption={Tasks Tablosu Semasi}]
CREATE TABLE IF NOT EXISTS tasks (
    id INTEGER PRIMARY KEY AUTOINCREMENT,
    location_id INTEGER REFERENCES locations(id),
    assigned_to INTEGER REFERENCES users(id),
    schedule_id INTEGER REFERENCES schedules(id),
    status TEXT CHECK(status IN ('pending', 'in_progress', 
        'completed', 'approved', 'rejected')) DEFAULT 'pending',
    priority TEXT CHECK(priority IN ('low', 'normal', 'high')) 
        DEFAULT 'normal',
    notes TEXT,
    due_date DATE,
    created_at DATETIME DEFAULT CURRENT_TIMESTAMP
);
\end{lstlisting}

% ============================================================================
% 5.4 KULLANIM DURUM DIYAGRAMI
% ============================================================================
\section{Kullanim Durum Diyagrami (Use-Case Diagram)}
\label{sec:use-case-diagram}

Sistemin kullanim durumlari, aktorlere gore Sekil 
\ref{fig:use-case}'de gosterilmektedir.

\begin{figure}[H]
\centering
\begin{tikzpicture}[
    usecase/.style={ellipse, draw=black, fill=cyan!20,
                    minimum width=2.8cm, minimum height=0.7cm,
                    text centered, font=\small},
    actor/.style={font=\small}
]
    % Sistem siniri
    \draw[dashed] (2,-6.5) rectangle (9,1);
    \node at (5.5,1.3) {\textbf{KTTYS}};
    
    % Aktorler
    \node[actor] at (0, 0) {Super Admin};
    \node[actor] at (0, -2) {Fakulte Admin};
    \node[actor] at (0, -4) {Supervisor};
    \node[actor] at (0, -6) {Personel};
    
    % Use cases
    \node[usecase] (uc1) at (5.5, 0) {Fakulte Yonetimi};
    \node[usecase] (uc2) at (5.5, -1.3) {Kullanici Yonetimi};
    \node[usecase] (uc3) at (5.5, -2.6) {Program Olusturma};
    \node[usecase] (uc4) at (5.5, -3.9) {Gorev Onaylama};
    \node[usecase] (uc5) at (5.5, -5.2) {Gorev Tamamlama};
    \node[usecase] (uc6) at (5.5, -6.2) {Sisteme Giris};
    
    % Iliskiler
    \draw (0.8, 0) -- (uc1);
    \draw (0.8, 0) -- (uc2);
    \draw (0.8, -2) -- (uc2);
    \draw (0.8, -2) -- (uc3);
    \draw (0.8, -4) -- (uc4);
    \draw (0.8, -6) -- (uc5);
    \draw (0.8, 0) -- (uc6);
    \draw (0.8, -2) -- (uc6);
    \draw (0.8, -4) -- (uc6);
    \draw (0.8, -6) -- (uc6);
    
\end{tikzpicture}
\caption{KTTYS Kullanim Durum Diyagrami}
\label{fig:use-case}
\end{figure}

Sekil \ref{fig:use-case}'deki kullanim durum diyagrami, sistemdeki 
dort farkli aktor tipinin gerceklestirebilecegi islemleri 
gostermektedir:

\begin{description}
    \item[Super Admin:] Sistem genelinde en yuksek yetkiye sahip 
    kullanicidir. Fakulte ekleme/silme ve tum kullanici yonetimi 
    islemlerini yapabilir.
    
    \item[Fakulte Admin:] Kendi fakultesi kapsaminda kullanici, 
    lokasyon, program ve gorev yonetimi yapabilir.
    
    \item[Supervisor:] Temizlik personelinin gorevlerini 
    denetler, tamamlanan gorevleri onaylar veya reddeder.
    
    \item[Personel:] Kendisine atanan temizlik gorevlerini 
    goruntler ve tamamlandi olarak isaretler.
\end{description}

% ============================================================================
% 5.5 AKIS DIYAGRAMI
% ============================================================================
\section{Akis Diyagrami (Flow Chart)}
\label{sec:akis-diyagrami}

Bu bolumde, sistemin temel is akislarinin diyagramlari 
sunulmaktadir.

\subsection{Kullanici Giris Akisi}
\label{subsec:giris-akisi}

\begin{figure}[H]
\centering
\begin{tikzpicture}[
    node distance=1.2cm,
    startstop/.style={rectangle, rounded corners, minimum width=2.5cm, 
        minimum height=0.7cm, text centered, draw=black, fill=red!30},
    process/.style={rectangle, minimum width=2.5cm, minimum height=0.7cm, 
        text centered, draw=black, fill=orange!30},
    decision/.style={diamond, minimum width=2cm, minimum height=0.7cm, 
        text centered, draw=black, fill=green!30, aspect=2},
    io/.style={trapezium, trapezium left angle=70, trapezium right angle=110,
        minimum width=2cm, minimum height=0.7cm, text centered, 
        draw=black, fill=blue!30},
    arrow/.style={thick,->,>=stealth}
]
    \node (start) [startstop] {Baslangic};
    \node (input) [io, below=of start] {Kullanici Adi/Sifre};
    \node (check1) [decision, below=of input] {Gecerli?};
    \node (query) [process, below=of check1] {Veritabani Sorgusu};
    \node (check2) [decision, below=of query] {Basarili?};
    \node (session) [process, below=of check2] {Oturum Olustur};
    \node (end) [startstop, below=of session] {Dashboard};
    \node (error) [process, right=2cm of check2] {Hata Mesaji};
    
    \draw [arrow] (start) -- (input);
    \draw [arrow] (input) -- (check1);
    \draw [arrow] (check1) -- node[right] {Evet} (query);
    \draw [arrow] (check1.east) -- ++(1,0) |- node[pos=0.25, above] {Hayir} (input.east);
    \draw [arrow] (query) -- (check2);
    \draw [arrow] (check2) -- node[right] {Evet} (session);
    \draw [arrow] (check2) -- node[above] {Hayir} (error);
    \draw [arrow] (error) |- (input);
    \draw [arrow] (session) -- (end);
\end{tikzpicture}
\caption{Kullanici Giris Akis Diyagrami}
\label{fig:giris-akisi}
\end{figure}

Sekil \ref{fig:giris-akisi}'daki akis diyagrami, kullanici giris 
surecinin adimlarini gostermektedir. Giris islemi su asamalardan 
olusmaktadir:

\begin{enumerate}
    \item Kullanici, kullanici adi ve sifresini girer.
    \item Sistem, alanlarin dolu olup olmadigini kontrol eder.
    \item Veritabanindan kullanici sorgulanir.
    \item bcrypt ile sifre dogrulamasi yapilir.
    \item Basarili giris sonrasi oturum olusturulur.
    \item Kullanici rolune uygun dashboard'a yonlendirilir.
\end{enumerate}

% ============================================================================
% 5.6 API TASARIMI
% ============================================================================
\section{API Tasarimi}
\label{sec:api-tasarimi}

KTTYS, RESTful prensiplerine uygun bir API sunmaktadir. Tablo 
\ref{tab:api-endpoints}'da temel API endpoint'leri listelenmistir.

\begin{longtable}{|p{2cm}|p{4.5cm}|p{5cm}|}
\hline
\textbf{Metod} & \textbf{Endpoint} & \textbf{Aciklama} \\
\hline
\endfirsthead
\hline
\textbf{Metod} & \textbf{Endpoint} & \textbf{Aciklama} \\
\hline
\endhead
\hline
\endfoot

POST & /api/auth/login & Kullanici girisi \\
\hline
POST & /api/auth/logout & Oturumu sonlandir \\
\hline
GET & /api/auth/me & Mevcut kullanici bilgisi \\
\hline
GET & /api/faculties & Tum fakulteleri listele \\
\hline
POST & /api/faculties & Yeni fakulte ekle \\
\hline
GET & /api/users & Kullanicilari listele \\
\hline
POST & /api/users & Yeni kullanici ekle \\
\hline
GET & /api/tasks & Gorevleri listele \\
\hline
POST & /api/tasks/:id/start & Goreve basla \\
\hline
POST & /api/tasks/:id/complete & Gorevi tamamla \\
\hline
POST & /api/tasks/:id/approve & Gorevi onayla \\
\hline
GET & /api/schedules & Programlari listele \\
\hline
POST & /api/schedules & Program ekle \\
\hline

\caption{KTTYS API Endpoint'leri}
\label{tab:api-endpoints}
\end{longtable}

% ============================================================================
% 5.7 KULLANICI ARAYUZU
% ============================================================================
\section{Kullanici Arayuzu}
\label{sec:kullanici-arayuzu}

KTTYS, modern ve kullanici dostu bir web arayuzu sunmaktadir. 
Arayuz, responsive tasarim prensipleriyle gelistirilmis olup, 
masaustu ve mobil cihazlarda optimal goruntuleme saglamaktadir.

\subsection{Arayuz Bilesenleri}
\label{subsec:arayuz-bilesenleri}

\begin{enumerate}
    \item \textbf{Giris Ekrani:} Kullanici adi ve sifre giris formu
    \item \textbf{Dashboard:} Role ozgu ozet bilgiler ve kisayollar
    \item \textbf{Fakulte Yonetimi:} Fakulte CRUD islemleri
    \item \textbf{Kullanici Yonetimi:} Personel listesi ve duzenleme
    \item \textbf{Lokasyon Yonetimi:} Bina/Kat/Oda hiyerarsisi
    \item \textbf{Program Yonetimi:} Haftalik program olusturma
    \item \textbf{Gorev Takip:} Gunluk gorev listesi ve durumlari
    \item \textbf{Raporlar:} Performans ve istatistik raporlari
\end{enumerate}

\subsection{Renk Paleti ve Tasarim}
\label{subsec:tasarim}

Sistem, modern ve profesyonel bir gorunum icin asagidaki 
renk paletini kullanmaktadir:

\begin{itemize}
    \item \textbf{Ana Renk:} \#1e40af (Koyu Mavi)
    \item \textbf{Basari:} \#16a34a (Yesil)
    \item \textbf{Uyari:} \#ea580c (Turuncu)
    \item \textbf{Hata:} \#dc2626 (Kirmizi)
    \item \textbf{Arkaplan:} \#f8fafc (Acik Gri)
\end{itemize}

Bu bolumde aciklanan sistem bilesenleri, KTTYS'nin tam 
fonksiyonel bir temizlik takip sistemi olarak calismasini 
saglamaktadir.
